\documentclass[10pt]{article}
\usepackage[margin=20mm]{geometry}
\usepackage{amsmath,amssymb,mathtools,fontspec,unicode-math,xurl,hyperref}
\IfFontExistsTF{DejaVu Serif}{\setmainfont{DejaVu Serif}}{\setmainfont{Latin Modern Roman}}
\setmathfont{latinmodern-math.otf}
\setlength{\parindent}{0pt}\setlength{\parskip}{6pt}\sloppy
\title{Scientific initial review of M25-0027}
\author{}
\date{25 October 2026}
\begin{document}
\maketitle
Abstract: The manifest-pinned 2022 manuscript has an incorrect numerical criterion for an affine moduli scheme. A genus-two, rank-one example satisfies the stated criterion but has additive stabilizers. A stronger criterion restores the auxiliary construction and is compatible with the later choice of a sufficiently large divisor. Separately, an explicit elliptic-curve calculation reproduces a core component-indexing gap that was corrected in a later author manuscript. The highest supported impact in the pinned version is Major, from that historical core proof gap; the unresolved threshold finding is Minor. These are findings about the identified versions, not findings about an unexamined journal-final PDF.\par
\section*{Identity, outcome, and scope}
The paper is Joakim Færgeman and Sam Raskin, Non-vanishing of geometric Whittaker coefficients for reductive groups. The assigned journal identifier is DOI 10.1090/jams/1051. The reviewed source is arXiv 2207.02955v1, submitted 6 July 2022; the PDF itself displays 8 July 2022. It has 64 pages. All page numbers below are one-based physical PDF pages, agreeing with the printed numbering. The source manifest classifies acquisition as version-limited; equivalence to the journal version is unestablished.\par
The highest supported impact in the pinned v1 manuscript is Major, owing to the independently reproduced historical component-indexing gap in Problem 2. The component classification is false for disconnected center and feeds directly into the main conservativity argument. This is a known historical issue reproduced; subsequent correction is credited, and it is not counted as a new unresolved defect. The unresolved numerical-threshold finding in Problem 1 is Minor: Definition 5.1.7.1 does not imply Lemma 5.1.7.2. A complete counterexample and a supported stronger hypothesis are given below. The later author’s complete replacement proof has not been independently verified; no journal-final equivalence is asserted.\par
All 64 pages were read as extracted text. Actual rendered pages 1, 5, 14, 17, 18, 27, 28, 35, 37, 39, 42–44, 46–48, 50, 52, and 53 were inspected visually. Selected exact TeX ranges were compared with the decisive displays. Reading all pages does not mean re-proving all results: the full microlocal index theorem, the entire critical-level localization construction, and the cited foundational results were not independently established.\par
\section*{Definitions, quantifiers, and dependency checks}
The manuscript works in characteristic zero, with a smooth geometrically connected projective curve \(X\) and a split reductive group \(G\). There is no global connected-center assumption (pp. 5, 14). The sheaf category in the main argument is the specified regular-holonomic/ind-constructible \(D\)-module context; its statements are not unrestricted assertions about arbitrary \(\ell\)-adic sheaves.\par
Equation reference: rootline. Fix a Borel \(B\) with unipotent radical \(N\) and Cartan \(T\). Write \(\check\rho\) for the half-sum of positive coroots, so that \(\langle\alpha_i,\check\rho\rangle=1\) for each simple root. After a choice of square root of \(\Omega_X\), the \(T\)-bundle used in Section 5 is \[\check\rho(\Omega_X):=(2\check\rho)(\Omega_X^{1/2}).\] For a coweight-valued divisor \(D\), the stack \(\operatorname{Bun}_N^{\Omega(-D)}\) is the fiber of \(\operatorname{Bun}_B\to\operatorname{Bun}_T\) over \(\check\rho(\Omega_X)(-D)\). Consequently its root subgroup for \(\alpha>0\) is the additive group of the line bundle \[
 L_\alpha=\Omega_X^{\langle\alpha,\check\rho\rangle}
                 \otimes\mathcal O_X(-\langle\alpha,D\rangle).\] The factor in this exponent is essential for the numerical check.\par
Theorems A and B concern the enhanced coefficient functor and conservativity; neither asserts full faithfulness. The first coefficient alone is not asserted to detect every tempered object before Hecke operations. The shifted coefficient functor \(\operatorname{coeff}[\dim\operatorname{Bun}_G]\) is the one whose exactness is used.\par
Several formal steps were checked conditionally on the cited inputs. In Section 7, exact baby-Whittaker averaging makes its kernel stable under truncations and subobjects, permitting the quotient \(t\)-structure. Right exactness of Hecke action is transferred through the conservative quotient functor; left exactness follows by adjunction using the dual representation. The fully faithful adjoint to this quotient is the left adjoint, as the manuscript states.\par
In Section 8.2 the identity \(\operatorname{coeff}\simeq\operatorname{coeff}_![-2d]\), \(d=\dim\operatorname{Bun}_G\), gives the stated Verdier-duality compatibility after shift by \(d\). In Lemma 10.1.1.1 the authors first reduce to a constructible object in the heart. Exactness then places its coefficient in a single degree, so the Euler-characteristic argument is legitimate in that reduction; it is not the invalid implication that an arbitrary complex with zero Euler characteristic vanishes. These checks do not remove Problems 1 and 2.\par
\section*{Problem 1: the large-divisor bound does not remove stabilizers}
\subsection*{Claim and location}
Equation reference: oldbound. Definition 5.1.7.1, equation (5.1.1), p. 28, declares \(D\) sufficiently large if \[
 \langle\alpha,\deg D\rangle>g-1
 \qquad(\alpha>0).\] Lemma 5.1.7.2 on the same page then asserts that \(\operatorname{Bun}_N^{\Omega(-D)}\) is an affine scheme. The exact source is `coeffs.tex`, lines 3126–3148. The reduction to additive groups does not prove the claim: a torsor stack can retain \(H^0\) as its automorphism group.\par
\subsection*{Counterexample with all parameters specified}
Work over \(k=\mathbb C\). Let \(X\) be the smooth projective completion of \[y^2=x^5-x,\] and let \(P\) be the point at infinity. The polynomial is square-free: at \(x=0\) its derivative is \(-1\), and at any nonzero root \(x^4=1\) its derivative is \(5x^4-1=4\). The degree-five hyperelliptic cover has six branch points, including infinity, so Riemann–Hurwitz gives genus \(g=2\). It is geometrically connected and smooth after its stated completion.\par
At \(P\), \(x\) and \(y\) have orders \(-2\) and \(-5\). Hence \(\eta=dx/y\) has order \(-3-(-5)=2\). It is nonzero and regular at the finite branch points, as is seen using \(y\) as a local parameter, and has no other zeroes. Thus \[\operatorname{div}(\eta)=2P,\qquad
 \Omega_X\simeq\mathcal O_X(2P).\] Choose the square root \(\mathcal O_X(P)\).\par
Equation reference: actualstack. Take the split reductive group \(G=PGL_2\) with its upper Borel. Its center is trivial, so this example is independent of Problem 2. There is a single positive root \(\alpha\), and \(\check\rho\) is an integral dominant coweight with \(\langle\alpha,\check\rho\rangle=1\). Set \[D=2\check\rho\cdot P.\] It is a dominant coweight-valued divisor and satisfies every inequality in [oldbound] (\#oldbound), since \(2>2-1\). However [rootline] (\#rootline) gives \[L_\alpha=\Omega_X(-2P)\simeq\mathcal O_X.\] Therefore the twisted unipotent group is the constant \(\mathbb G_a\), and its torsor stack has isomorphism classes \(H^1(X,\mathcal O_X)\) and automorphism group \(H^0(X,\mathcal O_X)\) at every object. Over \(\mathbb C\), after a linear splitting, it is \[
 \operatorname{Bun}_{\mathbb G_a}\simeq \mathbb A^2\times B\mathbb G_a.\] In particular the trivial torsor has automorphism group \(\mathbb G_a\ne1\). A scheme has trivial automorphism groups at its geometric points. Thus the actual stack is not a scheme, contradicting the lemma, and a fortiori is not an affine scheme.\par
\subsection*{Alternative readings and supported repair}
The object is defined as a fiber over a fixed \(T\)-bundle, not a coarse moduli space of isomorphism classes. Passing silently to \(H^1\) would change that definition and its sheaf operations. The choice of theta characteristic does not halve the root bundle: the displayed definition uses \(2\check\rho\) on \(\Omega_X^{1/2}\), and [rootline] (\#rootline) follows. Both points were checked in the actual p. 27 image.\par
Equation reference: newbound. A sufficient replacement, for the reductive groups under discussion, is \[
 \langle\alpha_i,\deg D\rangle>2g-2
 \quad\hbox{for every simple root }\alpha_i.\] Here is the proof of the auxiliary assertion under this new hypothesis. If \(\alpha=\sum_i n_i\alpha_i>0\), then \(\langle\alpha,\check\rho\rangle=\sum_i n_i\). Thus \[\deg L_\alpha
 =(2g-2)\sum_i n_i-\sum_i n_i
       \langle\alpha_i,\deg D\rangle<0.\] It follows that \(H^0(X,L_\alpha)=0\); also \(H^2(X,L_\alpha)=0\) because \(X\) is a curve. Use the \(T\)-stable central filtration of \(N\) by positive-root height. Its successive central quotients, after twisting, are sums of these additive line bundles. For a central extension by such a vector group, the obstruction to lifting a torsor lies in \(H^2\), the lifts form a torsor under \(H^1\), and relative automorphisms lie in \(H^0\). The first and last vanish. Inductively, therefore, the moduli of lifts is a torsor under a finite-dimensional vector group over the previous affine scheme. Such a torsor over an affine scheme is affine (indeed trivial after using vanishing of first quasi-coherent cohomology). Starting with the trivial group proves that \(\operatorname{Bun}_N^{\Omega(-D)}\) is an affine scheme under [newbound] (\#newbound). This is a strengthened hypothesis, not a proof of the original numerical claim.\par
\subsection*{Impact and calibration}
Step 1 of Theorem 8.3.0.1 (p. 43) uses this affineness to bound cohomological amplitude, so its invocation for every \(D\) satisfying [oldbound] (\#oldbound) is unsupported. Nevertheless Step 2 explicitly permits \(\check\lambda=2n\check\rho\) with \(n\gg0\) at a chosen point. Such a divisor satisfies [newbound] (\#newbound) as well. The same Steps 2 and 3 then transfer exactness to \(D=0\) and to general \(D\), conditional on the other stated inputs. The incorrect numerical threshold is therefore locally repairable without altering that main exactness statement or the allowed class of groups. This is the reason for the Minor calibration; mere ease of repair is not being used to dismiss a damaged core conclusion. No counterexample to the coefficient exactness theorem is claimed.\par
\section*{Problem 2: a corrected historical component-indexing gap}
\subsection*{Old claim and its precise failure}
Proposition 2.5.4.1(3), p. 17, asserts connectedness of \(\mathring{\operatorname{Nilp}}^{\check\lambda}\) whenever all \(\langle\alpha_i,\check\lambda\rangle+2g-2\) are nonnegative. The next paragraph declares the component-indexing map \(c\) injective. Section 2.5.6, p. 18, identifies the components of the everywhere regular locus with \(\operatorname{Hom}(\mathbb G_m,Z_G)\) and identifies the entire degree \((2-2g)\check\rho\) stratum with the Kostant component. The TeX confirms these statements at lines 1955–2002 and 2123–2148.\par
Let \(k=\mathbb C\), \(G=SL_2\), and let \(X\) be the elliptic curve \(y^2=x^3-x\) with origin \(O\) at infinity. Fix a nowhere-zero differential \(\eta\) and the trivial square root of \(\Omega_X\). Take \(\check\lambda=0\). All numerical hypotheses are satisfied since \(g=1\). This is a split reductive group and an allowed curve.\par
There are four distinct two-torsion line bundles \[L_P=\mathcal O_X(P-O),\qquad
 P\in\{O,(0,0),(1,0),(-1,0)\}.\] For the three finite points, the divisor of \(x-x(P)\) is \(2P-2O\), giving \(L_P^2\simeq\mathcal O_X\); the usual identification of \(X\) with \(\operatorname{Pic}^0(X)\) distinguishes them. Choose these square trivializations and define determinant-trivial rank-two bundles and Higgs fields \[E_P=L_P\oplus L_P,\qquad
 \phi_P=\begin{pmatrix}0&1\\0&0\end{pmatrix}\eta.\] Each field is trace-zero, squares to zero, and has rank one everywhere. Its unique kernel line is \(L_P\), of degree zero. Hence every point belongs to the asserted connected stratum.\par
Conversely, in this stratum a generically nonzero map \(L^{-1}\to L\otimes\Omega_X\) has degree zero; therefore it is an isomorphism. The kernel line consequently satisfies \(L^2\simeq
\mathcal O_X\). Its isomorphism class defines a map to the finite etale four-point space \(\operatorname{Pic}^0(X)[2]\). All four fibers are nonempty by the construction above and are open and closed. Thus the stratum cannot be connected.\par
One can count its components exactly. If \(M=\left(\begin{smallmatrix}a&b\\c&d\end{smallmatrix}\right)\) and \(f=\left(\begin{smallmatrix}0&1\\0&0\end{smallmatrix}\right)\), then \[Mf-fM=\begin{pmatrix}-c&a-d\\0&c\end{pmatrix}.\] With \(\det M=1\), centralizing \(f\) means \(c=0\), \(d=a\), \(a^2=1\). The centralizer is \(\mu_2\times\mathbb G_a\), using \(b/a\) as the additive coordinate. The everywhere regular stratum is its torsor stack. \(\operatorname{Bun}_{\mu_2}\) has four components, while \(\operatorname{Bun}_{\mathbb G_a}\simeq\mathbb A^1\times B\mathbb G_a\) is connected. There are exactly four components, whereas \(\operatorname{Hom}(\mathbb G_m,\mu_2)=0\) predicts one. This calculation verifies the discrepancy directly; no sheaf with an arbitrary prescribed singular support has been assumed to exist.\par
\subsection*{Old core dependency and counterevidence}
The old proof of Theorem 9.1.0.1 chooses an entire degree stratum inside \(\operatorname{SS}(\mathcal F)\) from the existence of one irreducible component meeting it (Section 9.3.2, p. 46, source lines 5077–5086). On p. 47 it then puts a particular canonical Higgs point into \(\operatorname{SS}(\mathcal F)\) solely from membership in that stratum. The example shows why the degree label alone cannot justify these steps. Its four components have the same degree, discrepancy zero, and first Chern class zero.\par
The chain continues through Corollary 10.1.1.2, the enhanced-coefficient conservativity statement, and Theorem 10.3.3.1 to Theorems A and B. The invalid classification is therefore a Major historical proof defect. This does not establish that the main conservativity theorem is false, nor that the four Higgs points themselves supply a counterexample to a statement about all sheaves.\par
\subsection*{Later-version attribution}
The author-hosted manuscript [FRnew], dated 9 May 2025, changes Proposition 2.5.4.1 to account for disconnected center and central torsors, restricts the Kostant-stratum equality, and incorporates central translations into Section 9 and the subsequent conservativity argument. It also includes an \(SL_2\) example. These changes corroborate the historical issue. The replacement global proof has not been fully audited here. Its Lemma 5.1.7.2 retains the numerical bound considered in Problem 1. This comparison does not establish equivalence of either author manuscript to the journal-final PDF.\par
\section*{Remaining dependencies and rejected overclaims}
Section 8 cites a result of Kevin Lin as forthcoming at the time of writing. A primary manuscript is now available as arXiv 2211.05282v1 [Lin]; its introduction and stated relationship between Poincare series and duality were consulted. This is not currently an absent source, but its full proof and all normalization identifications have not been independently checked. The cited alternative topological argument was not audited.\par
Section 11 distinguishes irreducible from very irreducible local systems; the latter has automorphism group exactly \(Z_{\check G}\). The finite-group example on p. 52 is compatible with this distinction. For semisimple \(G\), the proof bounds the length of the normalized eigensheaf by \(|\pi_1^{\mathrm{alg}}(G)|\) and obtains the reverse bound from nonzero restrictions to every component. Conditional on exact conservative coefficients, equality forces length one on each component. For general reductive \(G\), the central coweight quotient addresses the infinite-component issue. These are conditional formal checks, not a reconstruction of the full existence theorem.\par
The paper itself acknowledges a general-parabolic extension of the Eisenstein-series input as folklore, with incomplete literature proofs (Remark 11.1.4.2). Appendix A invokes substantial oper and localization machinery. These remain external-proof limitations; they are not converted into additional established defects. The semisimplicity statement for arbitrary irreducible local systems is presented as expected in Remark 11.1.4.3 and is not accepted here as an unconditional theorem. No claims about unperformed experiments or author-code execution are made.\par
\section*{Conclusion}
Item | Judgment | Repair or remaining work\par
Large-divisor criterion | Minor, proved counterexample | Stronger simple-root bound verified; later use can select it\par
Old component indexing | Major in pinned v1; known historical issue | Later correction credited; complete replacement proof unaudited\par
Other selected formal deductions | No additional issue found | Conditional on listed geometric and categorical inputs\par
Journal-final attribution | Not established | Compare the actual final version before transferring findings\par
Acceptance and adoption | Human acceptance pending | No author adoption requested or asserted\par
\section*{Supporting checks}
All seven manifest-listed input files matched their recorded SHA-256 hashes. The pinned PDF hash is\par
> 9336d1df58a7d64f5d0863cc1dfd97b7f168
> 66454ddeab6c63017793291c5604\par
and the `coeffs.tex` hash is\par
> cd975980ae76b2e9946afadd58ede885352
> f89bfdc44a78b3f3a85af72ca76c3.\par
The source TeX ranges read directly were 1940–2005, 2110–2160, 3110–3155, and 5060–5140. The rest of the page coverage is based on the actual pinned PDF text, with visual checks listed above.\par
The reviewer-authored utility `verification/review\_io.py` validated input hashes, called local `pdfinfo` and `pdftotext`, and rendered selected original pages with `pdftoppm`. For example:\par
{\ttfamily\footnotesize python -B verification/review\_io.py inventory 0027\\
python -B verification/review\_io.py pages 0027 20 23\\
python -B verification/review\_io.py render 0027 17,18,27,28}\par
The actual commands used absolute paths within the authorized staging root and rendered the full page list recorded in the coverage ledger. These are source-processing commands, not execution of author programs. The mathematical counterexamples and the repair argument above were checked analytically; no numerical experiment or symbolic-software validation is claimed.\par
Evidence files include `input\_validation.json`, `source\_pages.json`, the actual `source\_page\_NNN.png` files, `visual\_inspection.json`, `source\_fragments.json`, and `compilation\_status.json`. The first rendered page produced a Poppler Symbol-font warning; the actual displayed page was legible and inspected. Rendering success alone was never treated as visual inspection.\par
\section*{Supporting checks}
The authorized followup uses the existing Tectonic 0.17.0 runtime. Exact-source build logs, input and PDF hashes, and report-page inspection records are retained in the accompanying provenance and followup evidence. No compiler was installed by this conversation. Original-source PDF inspection remains recorded separately. Coverage, findings, provenance, and the current report files are bound in delivery state after completion of the report-page checks.\par
No independent backend-route evidence or equivalence of names is asserted, and no dispatcher tab was inspected or taken over. The current assistant performed the scientific review, source checks, and actual page inspection. Human scientific acceptance remains pending. No author code, contact, installation, cloud upload, or writes outside this task staging tree were used.\par
9 J. Færgeman and S. Raskin, Non-vanishing of geometric Whittaker coefficients for reductive groups, arXiv:2207.02955v1 (6 July 2022), manifest-pinned 64-page source. Assigned journal DOI: 10.1090/jams/1051.\par
J. Færgeman and S. Raskin, same title, author-hosted 77-page manuscript, last updated 9 May 2025, https://www.samraskin.net/coeffs.pdf. Selected version-comparison passages consulted 25 October 2026 and rechecked against retained PDF bytes on 26 October 2026.\par
K. Lin, Poincare series and miraculous duality, arXiv:2211.05282v1 (10 November 2022), https://arxiv.org/abs/2211.05282v1. Introduction and dependency description consulted; full proof not audited.\par
\end{document}
